% Kodiranje iste poruke: JSON naspram Protocol Buffers
\documentclass[border=4pt]{standalone}
\usepackage[T1]{fontenc}\usepackage[utf8]{inputenc}
\usepackage[english,serbian]{babel}
\usepackage{tikz}\usetikzlibrary{arrows.meta,positioning,fit,backgrounds}
\tikzset{
  okvir/.style={draw=black,line width=0.5pt,inner sep=6pt,align=left},
  o/.style={font=\scriptsize,align=center},
  s/.style={-{Latex[length=1.8mm,width=1.5mm]},line width=0.4pt,black},
}
\begin{document}
\begin{tikzpicture}

  \node[font=\small\itshape,anchor=south] at (0,1.85)
       {ista poruka: \texttt{Artist \{ id = 1, name = "Monet" \}}};

  % ---------- JSON ----------
  \node[okvir,anchor=north,minimum width=6.6cm] (json) at (-3.6,1.2)
       {\ttfamily\footnotesize \{"id":1,"name":"Monet"\}};
  \node[o,anchor=north,text width=6.4cm] at (-3.6,-0.15)
       {\textbf{JSON} --- 23 bajta\\[3pt]
        nazivi polja \texttt{id} i \texttt{name}\\
        prenose se uz svaku vrednost};

  % ---------- PROTOBUF ----------
  \node[okvir,anchor=north,minimum width=6.6cm] (pb) at (3.6,1.2)
       {\ttfamily\footnotesize 08 01 12 05 4D 6F 6E 65 74};
  \node[o,anchor=north,text width=6.4cm] at (3.6,-0.15)
       {\textbf{Protocol Buffers} --- 9 bajtova\\[3pt]
        prenose se samo oznake polja\\
        \texttt{1} i \texttt{2} iz opisa poruke};

  % ---------- objasnjenje bajtova ----------
  \node[o,anchor=north,text width=13.6cm] at (0,-1.85)
       {\texttt{08}\,=\,oznaka 1, ceo broj \quad \texttt{01}\,=\,vrednost \quad
        \texttt{12}\,=\,oznaka 2, niska \quad \texttt{05}\,=\,dužina \quad
        \texttt{4D\,6F\,6E\,65\,74}\,=\,,,Monet``};

  \node[o,anchor=north,font=\small\itshape,text width=13cm] at (0,-2.75)
       {ušteda potiče od izostavljanja naziva polja, pa je najveća
        kod poruka sa mnogo malih polja, a najmanja kod poruka
        u kojima preovlađuje dug tekst};

\end{tikzpicture}
\end{document}
